\section{What the existing calculations contribute}\label{sec:constraints}

The previous calculations provide candidate ingredients and
restrictions on a realization. They do not yet supply the effective
transfer of Hypothesis~\ref{hyp:realization}.

\subsection{A spectral ingredient, with a fixed subtraction sign}

A free three-dimensional defect scalar in logarithmic radial
coordinates has angular rates $\ell+1/2$. Even Poisson smearing
selects $\ell=2n$ and gives, for $0\leq\rho<1$,
\begin{equation}
 C_\rho(u)=\sum_{n\geq0}\rho^{4n}e^{-a_n|u|}
       =\frac{e^{-|u|/2}}{1-\rho^4e^{-2|u|}}.
\end{equation}
This is a positive covariance tending to $\nGamma(|u|)$ off the
diagonal. If $T_\rho$ denotes convolution by $C_\rho$ and
$M_\rho=\widehat C_\rho(0)$, the positive difference form is
$M_\rho I-T_\rho$, whose limit is $\mathcal E_\Gamma$.
The opposite subtraction has the wrong high-frequency sign, and
a finite contact constant cannot repair it. The auxiliary Gaussian
mode model also produces the gamma partition ratio, but neither
a physical projection to the required multiplicities nor its
identification as a boundary transfer is established
\cite[Section~S2]{Supplement}.

These results explain why defect matter is relevant and fix a
sign that a candidate must respect. They supply neither $w_0$
by themselves nor the pole and arithmetic terms in
\eqref{eq:generator}. The auxiliary completion in
Section~\ref{sec:fixed-window} includes both pole factors by a
prescribed change of modes; their physical origin remains open. The exact positive Beta/comb factorization
likewise retains an essential rational subtraction
\cite[Section~S1.4]{Supplement}.

\subsection{Contour freedom and endpoint compatibility}

The inherited common-reference semicircles become vertical slits
after inversion and have constant Loewner driving. Their rigidity
holds for one fixed scalar prescription and one specified bulk
charge. A tangent-dependent scalar coupling changes that
restriction: with the full chiral D3--D5 projector, an oriented
real tangent isometry admits a common complex charge space of
dimension two in a normal plane and dimension one across all
defect directions. These are exact statements in the stated
Euclidean Clifford module
\cite[Sections~S3 and S5]{Supplement}.

Contour freedom does not complete the observable. In this
constant-map Poincare-charge ansatz, ordinary scalar endpoints
annihilated by the same charge have zero mutual free contraction.
Their bilinear also has nonzero residual R charge, giving a
vanishing expectation in an invariant vacuum. Pairing an endpoint
with its physical adjoint restores a nonzero contraction but changes
the supersymmetry conditions
\cite[Section~S6]{Supplement}. Baker's original conformal
semicircle is a nonzero control outside that restricted ansatz,
not an excluded construction.

A candidate should therefore be tested for a nonzero normalized
endpoint pairing before trying to extract its arithmetic generator.
A shared line supercharge alone does not ensure such a pairing.

\subsection{A nontrivial response, with an unfinished physical sum}

Ordinary time reflection gives a consistent classical adjoint and
a positive free endpoint kernel. A smooth polynomial bulge can
join its reflected branch without a fixed-angle cusp. In the
direct Feynman-gauge Gaussian bulk-exchange sector, the stripped
self-response is
\begin{equation}
 I(\eta)=\left(\tfrac32-2\log2\right)\eta^2+O(\eta^4),
 \qquad \tfrac32-2\log2>0.
\end{equation}
The construction and remainder estimate are in
\cite[Sections~S7--S8]{Supplement}.
Endpoint-to-line exchange, defect interactions, self-energies and
counterterms are still needed for the full response. Thus the
coefficient is a concrete insertion test, not an effective
arithmetic generator or an all-order protection theorem.

The distinction between supersymmetry and a positive norm is
visible even abstractly: $Q\psi=0$ and $\delta\psi=Q\lambda$
give $\delta\norm\psi^2=2\Ree\ip{Q^\dagger\psi}{\lambda}$,
which need not vanish. A successful reduction must handle both
the chosen charge and the physical adjoint.
